\section{Глава~\ref{sec:motorlearning}. Обучение движениям}

Правило наименьших квадратов и выгода отдельного провода ошибки --- теоремы; устройство схемы с проводом ошибки, ослабление при совпадении, подстройка следа под задержку, последействие и спонтанный возврат навыка измерены; то, что вся схема работает как обучение с учителем и строит внутреннюю модель, --- ведущая гипотеза; числа двухпроцессной модели --- расчёт по модели книги.

{\footnotesize
\setlength{\tabcolsep}{3pt}\renewcommand{\arraystretch}{1.15}
\begin{longtable}{>{\raggedright}p{0.5cm}>{\raggedright}p{4.0cm}>{\raggedright}p{5.6cm}>{\raggedright}p{2.3cm}>{\raggedright\arraybackslash}p{1.7cm}}
\toprule
№ & Утверждение & Опыт и что измерено & Источник & Статус \\
\midrule\endfirsthead
\toprule
№ & Утверждение & Опыт и что измерено & Источник & Статус \\
\midrule\endhead
1 & Правило~\eqref{eq:lms} в среднем идёт по градиенту квадрата ошибки & Результат теории адаптивных фильтров: поправка, пропорциональная входу и ошибке выхода, есть стохастический спуск по градиенту среднего квадрата ошибки. & \cite{WidrowHoff1960} & теория \\
2 & С проводом ошибки на каждый выход скорость не зависит от числа выходов; с одним общим числом --- падает с числом шумящих нейронов (утверждение~\ref{prop:teachers}) & Кривые обучения линейной сети: обучение по случайным возмущениям нейронов медленнее точного градиента во столько раз, сколько возмущаемых нейронов. & \cite{Werfel2005} & теория \\
3 & Схема с проводом ошибки работает как обучение с учителем (утверждение~\ref{prop:errorwire}) & Теории, выведенные из строения схемы. Опыты строк 7--10 согласуются с ними; что вся схема работает именно так, не доказано. & \cite{Marr1969, Albus1971} & гипотеза \\
4 & Расширитель: около $7\cdot10^{10}$ мелких \nt{GLUT}-нейронов, больше, чем во всём остальном мозге & Подсчёт ядер во взвеси растворённой ткани: в мозжечке человека $69\cdot10^9$ нейронов из $86\cdot10^9$ во всём мозге. Стереология: $101\cdot10^9$ мелких нейронов в мозжечке пожилых мужчин. & \cite{Azevedo2009, Andersen1992cereb} & измерено \\
5 & Подъём размерности входа делает нужную зависимость линейной & Теорема о числе разбиений: взвешенная сумма $n$ входов с порогом реализует произвольную разметку примерно $2n$ векторов общего положения. Что мозжечок пользуется именно этим, --- часть гипотезы строки~3. & \cite{Cover1965} & теория \\
6 & Около $3\cdot10^7$ выходных \nt{GABA}-нейронов, у каждого порядка $10^5$ входов & Стереология мозжечка человека: $30.5\cdot10^6$ выходных нейронов. Электронная микроскопия, крыса: около $1.75\cdot10^5$ входов от расширителя на один выходной нейрон. Тормозной медиатор выходных нейронов --- в обзоре. & \cite{Andersen1992cereb, NapperHarvey1988, Ito2001} & измерено \\
7 & Ровно один провод ошибки на выходной нейрон, около раза в секунду, каждый раз мощная вспышка & Обзор записей: каждый выходной нейрон получает один лазящий \nt{GLUT}-вход, который вызывает сложный импульс с частотой порядка $1$~Гц. & \cite{Ito2001} & измерено \\
8 & Вероятность вспышки зависит от направления ошибки движения & Обезьяна, глазодвигательный отдел мозжечка, адаптация саккад: направление зрительной ошибки задаёт вероятность сложного импульса, величина --- его время; вспышка меняет простые импульсы в следующей попытке и сдвигает движение вдоль предпочтительной ошибки клетки. & \cite{Herzfeld2018} & измерено \\
9 & Совпадение вспышки ошибки с активным входом ослабляет вход ($-\eta e_i x_j$) & Обзор опытов: совместная стимуляция входа от расширителя и провода ошибки вызывает долговременное ослабление этого входа; стимуляция одного входа --- нет. & \cite{Ito2001} & измерено \\
10 & Длина следа подстроена под задержку ошибки: при задержке $100\ms$ сильнее всего ослабляется вход, бывший активным за $100\ms$ & Срезы и живые мыши: в отделе, отвечающем за глазодвигательное обучение, ослабление узко настроено на интервал около $120\ms$ между входом и ошибкой --- столько идёт сигнал ошибки в эту задачу; в другом отделе разные клетки настроены на разные интервалы. & \cite{Suvrathan2016} & измерено \\
11 & Схема --- адаптивный фильтр~\eqref{eq:adaptivefilter}, выучивающий внутреннюю модель тела & Модель, выведенная из строения схемы. Прямого опыта, где выученный фильтр измерен как передаточная функция тела, нет. & \cite{Fujita1982} & гипотеза \\
12 & Предсказание вычитает следствия собственных движений, поэтому себя не пощекотать & Люди: собственное касание кажется менее щекотным, чем то же чужое; в томографе соматосенсорная кора отвечает на собственное касание слабее, а мозжечок меньше активен, когда движение касается кожи. Объяснение через вычитание предсказания --- гипотеза. & \cite{Blakemore1998} & гипотеза \\
13 & При внешней силе промах уменьшается, а после её снятия рука промахивается в другую сторону & Люди тянутся рукоятью робота в силовом поле: траектории возвращаются к прямым; после снятия поля последействие --- почти зеркальное отражение первых искажений; оно переносится и в не тренированные части пространства. & \cite{ShadmehrMussaIvaldi1994} & измерено \\
14 & Учатся два процесса~\eqref{eq:tworate}; после обратного возмущения навык возвращается сам & Люди, силовое поле, затем короткое обратное поле и попытки с зажатой ошибкой: поправка сама возвращается к старой; модель из двух процессов это воспроизводит, модель из одного --- нет. & \cite{Smith2006} & измерено \\
15 & Числа рис.~\ref{fig:tworate}: $0.84$ (медленный $0.63$) за $200$ движений; $6$ обратных; быстрый $-0.53$, медленный $0.46$; возврат до $0.37$ за $40$ движений & Расчёт по \texttt{models/\allowbreak motor\_\allowbreak learning.py}, $A_f=0.92$, $B_f=0.1$, $A_s=0.996$, $B_s=0.02$: $0.840$ ($0.634$ и $0.205$), $6$ движений, $-0.531$ и $0.457$, пик $0.371$ на $40$-м движении. & \texttt{models/\allowbreak motor\_\allowbreak learning.py} & модель \\
16 & Два процесса нужны, потому что тело меняется с разными скоростями & Теория: оптимальная оценка при помехах с несколькими временами жизни даёт несколько фильтров с разными постоянными времени и объясняет спонтанный возврат и ускоренное повторное обучение. & \cite{Kording2007} & гипотеза \\
\bottomrule
\end{longtable}}
